%% Font size %%
\documentclass[11pt]{article}

%% Load the custom package
\usepackage{Mathdoc}

%% Numéro de séquence %% Titre de la séquence %%
\renewcommand{\centerhead}{S14- Probabilités - Exercices Divers}

%% Spacing commands %%
\renewcommand{\baselinestretch}{1} \setlength{\parindent}{0pt}

%%\pagestyle{empty}

\begin{document}

\begin{exercice}[1][Reconnaître une expérience aléatoire]

Pour chaque situation, dire s'il s'agit ou non d'une expérience aléatoire.

\begin{enumerate}
    \item Lancer un dé à 6 faces.
    \item Mesurer la longueur d'une table avec une règle.
    \item Tirer une carte dans un jeu.
    \item Choisir au hasard une boule dans un sac.
    \item Calculer $15+7$.
\end{enumerate}

\end{exercice}


\begin{exercice}[1][Décrire les issues possibles]

On lance une pièce de monnaie.

\begin{enumerate}
    \item Donner toutes les issues possibles.
    \item Combien y a-t-il d'issues possibles ?
\end{enumerate}

\end{exercice}


\begin{exercice}[1][Issues d'un dé]

On lance un dé à 6 faces numérotées de 1 à 6.

\begin{enumerate}
    \item Donner toutes les issues possibles.
    \item Donner l'événement :
    \begin{enumerate}
        \item « obtenir un nombre pair » ;
        \item « obtenir un nombre supérieur à 4 ».
    \end{enumerate}
\end{enumerate}

\end{exercice}


\begin{exercice}[1][Événements]

On tire une boule dans un sac contenant des boules rouges, bleues et vertes.

Associer chaque événement à ses issues possibles.

\begin{enumerate}
    \item « Obtenir une boule rouge »
    \item « Obtenir une boule qui n'est pas verte »
    \item « Obtenir une boule bleue ou verte »
\end{enumerate}

\end{exercice}

\newpage


\begin{exercice}[2][Calculer une probabilité]

On lance un dé à 6 faces.

Calculer la probabilité des événements suivants.

\begin{enumerate}
    \item Obtenir un 3.
    \item Obtenir un nombre pair.
    \item Obtenir un nombre inférieur à 5.
    \item Obtenir un nombre strictement supérieur à 2.
\end{enumerate}

\end{exercice}


\begin{exercice}[2][Boules dans un sac]

Un sac contient 5 boules rouges, 3 boules bleues et 2 boules vertes.

On tire une boule au hasard.

\begin{enumerate}
    \item Combien y a-t-il de boules au total ?
    \item Calculer la probabilité d'obtenir :
    \begin{enumerate}
        \item une boule rouge ;
        \item une boule bleue ;
        \item une boule verte.
    \end{enumerate}
\end{enumerate}

\end{exercice}


\begin{exercice}[2][Cartes numérotées]

On place dans une boîte des cartes numérotées de 1 à 10.

On tire une carte au hasard.

Calculer la probabilité des événements suivants.

\begin{enumerate}
    \item Obtenir un nombre pair.
    \item Obtenir un multiple de 3.
    \item Obtenir un nombre supérieur ou égal à 8.
    \item Obtenir un nombre inférieur à 11.
\end{enumerate}

\end{exercice}


\begin{exercice}[2][Roue de loterie]

Une roue est partagée en 8 parts identiques :
3 rouges, 2 bleues, 2 vertes et 1 jaune.

On fait tourner la roue.

\begin{enumerate}
    \item Calculer la probabilité d'obtenir la couleur rouge.
    \item Calculer la probabilité d'obtenir la couleur verte.
    \item Calculer la probabilité de ne pas obtenir la couleur jaune.
\end{enumerate}

\end{exercice}


\begin{exercice}[2][Vrai ou faux]

Dire si les affirmations suivantes sont vraies ou fausses.

\begin{enumerate}
    \item Une probabilité peut être supérieure à 1.
    \item Une probabilité peut être égale à 0.
    \item Une probabilité peut être négative.
    \item La somme des probabilités de toutes les issues vaut 1.
\end{enumerate}

\end{exercice}


\newpage


\begin{exercice}[2][Événement contraire]

Dans une boîte, il y a 9 boules blanches et 6 boules noires.

On tire une boule au hasard.

\begin{enumerate}
    \item Calculer la probabilité d'obtenir une boule noire.
    \item En déduire la probabilité d'obtenir une boule blanche.
\end{enumerate}

\end{exercice}


\begin{exercice}[2][Compléter avec l'événement contraire]

On lance un dé à 6 faces.

Compléter.

\begin{enumerate}
    \item L'événement contraire de « obtenir un nombre pair » est :
    \dotfill

    \item L'événement contraire de « obtenir un nombre inférieur à 5 » est :
    \dotfill

    \item L'événement contraire de « obtenir un 6 » est :
    \dotfill
\end{enumerate}

\end{exercice}


\begin{exercice}[2][Événement certain ou impossible]

Pour chaque événement, préciser s'il est :

\begin{center}
certain \hspace{1cm} impossible \hspace{1cm} possible
\end{center}

\begin{enumerate}
    \item Obtenir un nombre inférieur à 7 avec un dé à 6 faces.
    \item Obtenir un 9 avec un dé à 6 faces.
    \item Obtenir un nombre pair avec un dé à 6 faces.
    \item Tirer une boule rouge dans un sac qui ne contient que des boules rouges.
\end{enumerate}

\end{exercice}


\begin{exercice}[3][Problème — Tombola]

Lors d'une tombola, 200 tickets sont vendus.

Parmi eux :

\begin{itemize}
    \item 1 ticket permet de gagner un vélo ;
    \item 4 tickets permettent de gagner un ballon ;
    \item 20 tickets permettent de gagner une peluche.
\end{itemize}

On choisit un ticket au hasard.

\begin{enumerate}
    \item Quelle est la probabilité de gagner le vélo ?
    \item Quelle est la probabilité de gagner un ballon ?
    \item Quelle est la probabilité de gagner un cadeau ?
    \item Quelle est la probabilité de ne rien gagner ?
\end{enumerate}

\end{exercice}


\begin{exercice}[3][Problème — Météo]

Un professeur note la météo pendant 30 jours.

\begin{center}
\begin{tabular}{|c|c|}
\hline
Temps & Nombre de jours \\
\hline
Soleil & 12 \\
Nuages & 10 \\
Pluie & 8 \\
\hline
\end{tabular}
\end{center}

On choisit un jour au hasard.

\begin{enumerate}
    \item Calculer la probabilité qu'il ait fait soleil.
    \item Calculer la probabilité qu'il ait plu.
    \item Calculer la probabilité qu'il n'ait pas plu.
\end{enumerate}

\end{exercice}


\newpage


\begin{exercice}[3][Tableau de probabilités]

Une urne contient des jetons rouges, bleus et verts.

Compléter le tableau.

\begin{center}
\begin{tabular}{|c|c|c|}
\hline
Couleur & Nombre de jetons & Probabilité \\
\hline
Rouge & 4 & \\
\hline
Bleu & 3 & \\
\hline
Vert & 5 & \\
\hline
Total & & 1 \\
\hline
\end{tabular}
\end{center}

\end{exercice}


\begin{exercice}[3][Comparer des probabilités]

Dans un premier sac, il y a 2 boules rouges et 3 boules bleues.

Dans un deuxième sac, il y a 5 boules rouges et 10 boules bleues.

\begin{enumerate}
    \item Calculer la probabilité d'obtenir une boule rouge dans chaque sac.
    \item Dans quel sac a-t-on le plus de chances d'obtenir une boule rouge ?
\end{enumerate}

\end{exercice}


\begin{exercice}[3][Chercher une quantité]

Dans un sac contenant 20 boules, la probabilité d'obtenir une boule rouge est $\dfrac{7}{20}$.

\begin{enumerate}
    \item Combien y a-t-il de boules rouges ?
    \item Combien y a-t-il de boules qui ne sont pas rouges ?
\end{enumerate}

\end{exercice}


\begin{exercice}[3][Jeu télévisé]

Dans un jeu, un candidat choisit au hasard une enveloppe parmi 50.

\begin{itemize}
    \item 5 enveloppes contiennent 100 € ;
    \item 15 enveloppes contiennent 20 € ;
    \item les autres enveloppes sont vides.
\end{itemize}

\begin{enumerate}
    \item Quelle est la probabilité de gagner 100 € ?
    \item Quelle est la probabilité de gagner de l'argent ?
    \item Quelle est la probabilité de ne rien gagner ?
\end{enumerate}

\end{exercice}


\begin{exercice}[4][Exercice de recherche]

Un jeu consiste à lancer un dé à 6 faces.

\begin{itemize}
    \item Si le joueur obtient 1 ou 2, il gagne 2 points.
    \item S'il obtient 3, 4 ou 5, il gagne 1 point.
    \item S'il obtient 6, il perd 1 point.
\end{itemize}

\begin{enumerate}
    \item Quelle est la probabilité de gagner 2 points ?
    \item Quelle est la probabilité de gagner au moins 1 point ?
    \item Quelle est la probabilité de perdre 1 point ?
\end{enumerate}

\end{exercice}


\begin{exercice}[4][Créer une expérience aléatoire]

Inventer une expérience aléatoire.

Puis :

\begin{enumerate}
    \item donner toutes les issues possibles ;
    \item proposer deux événements ;
    \item calculer la probabilité d'un des événements.
\end{enumerate}

\end{exercice}

\end{document}
